\section{Discussion: Generalization and Boundaries}
\label{sec:discussion}

The framework-first view separates the data-construction method from the rotation case study. The broader idea is a reusable data engine for object-level editing tasks that require controlled changes and preserved invariants. The specific yaw contract does not transfer automatically to other tasks, but the workflow structure can: staged artifact construction, review, bounded feedback actions, downstream validation, and explicit verdicts.

This generalization must be stated cautiously. The workflow does not automatically solve resizing, repositioning, material editing, insertion, pose editing, or real-data collection. Each task must define its own change/invariance contract, artifact schema, quality signals, action space, and acceptance logic. The framework provides a disciplined way to express those requirements and inspect whether the resulting dataset satisfies them.

The next extension path is mixed or real-data construction. A future version could use VLM feedback to review collected real samples, identify missing coverage or quality issues, and convert downstream weaknesses into either synthetic augmentation plans or real-data collection plans. The same quality-gating logic could then apply across data sources. This report does not claim that the system already solves real-data construction; it identifies a plausible extension of the same closed-loop abstraction.

The practical lesson is that quality gates belong inside the data engine. If low-quality samples can enter training merely because they were generated, scaling becomes unsafe. The next version of the workflow should treat render quality, asset integrity, view-instruction consistency, and category-aware rendering priors as first-class acceptance conditions.
